% =============================================================================
% DS – Chapitre 5 : Résolution d'équations du 1er degré
% 2de Bac Pro | 3 versions : Socle – Standard – Approfondissement
% Compilable avec pdflatex (Overleaf compatible)
% =============================================================================
\documentclass[a4paper,11pt]{article}

% --- Encodage et langue ---
\usepackage[utf8]{inputenc}
\usepackage[T1]{fontenc}
\usepackage[french]{babel}

% --- Mise en page ---
\usepackage[margin=2cm]{geometry}
\usepackage{parskip}
\usepackage{enumitem}
\usepackage{multicol}

% --- Maths ---
\usepackage{amsmath,amssymb}

% --- Tableaux ---
\usepackage{booktabs}
\usepackage{array}

% --- Couleurs et boîtes ---
\usepackage[dvipsnames]{xcolor}
\usepackage[most]{tcolorbox}

% --- En-têtes ---
\usepackage{fancyhdr}
\usepackage{lastpage}

% --- Divers ---
\usepackage{tikz}
\usepackage{eurosym}

% =============================================================================
% Couleurs des compétences
% =============================================================================
\definecolor{appblue}{HTML}{3B82F6}
\definecolor{anapurple}{HTML}{8B5CF6}
\definecolor{reagreen}{HTML}{22C55E}
\definecolor{valorange}{HTML}{F59E0B}
\definecolor{comyellow}{HTML}{EAB308}
\definecolor{corrgray}{HTML}{F3F4F6}
\definecolor{corrgreen}{HTML}{166534}
\definecolor{soclecol}{HTML}{22C55E}
\definecolor{stdcol}{HTML}{3B82F6}
\definecolor{approcol}{HTML}{8B5CF6}

% =============================================================================
% Badges de compétences
% =============================================================================
\newtcbox{\badgeAPP}{on line, colback=appblue!15, colframe=appblue!60,
  fontupper=\scriptsize\bfseries\color{appblue}, boxrule=0.5pt,
  arc=3pt, boxsep=1pt, left=3pt, right=3pt, top=1pt, bottom=1pt}
\newtcbox{\badgeANA}{on line, colback=anapurple!15, colframe=anapurple!60,
  fontupper=\scriptsize\bfseries\color{anapurple}, boxrule=0.5pt,
  arc=3pt, boxsep=1pt, left=3pt, right=3pt, top=1pt, bottom=1pt}
\newtcbox{\badgeREA}{on line, colback=reagreen!15, colframe=reagreen!60,
  fontupper=\scriptsize\bfseries\color{reagreen}, boxrule=0.5pt,
  arc=3pt, boxsep=1pt, left=3pt, right=3pt, top=1pt, bottom=1pt}
\newtcbox{\badgeVAL}{on line, colback=valorange!15, colframe=valorange!60,
  fontupper=\scriptsize\bfseries\color{valorange}, boxrule=0.5pt,
  arc=3pt, boxsep=1pt, left=3pt, right=3pt, top=1pt, bottom=1pt}
\newtcbox{\badgeCOM}{on line, colback=comyellow!15, colframe=comyellow!60,
  fontupper=\scriptsize\bfseries\color{comyellow!80!black}, boxrule=0.5pt,
  arc=3pt, boxsep=1pt, left=3pt, right=3pt, top=1pt, bottom=1pt}

% =============================================================================
% Boîte de correction
% =============================================================================
\newtcolorbox{correction}{
  colback=corrgray,
  colframe=corrgreen!70!black,
  coltitle=white,
  fonttitle=\bfseries,
  title=Correction,
  arc=4pt,
  boxrule=0.8pt,
  left=6pt, right=6pt, top=4pt, bottom=4pt,
  before skip=8pt, after skip=8pt
}

% =============================================================================
% Badge de version
% =============================================================================
\newcommand{\versionsocle}{%
  \begin{tcolorbox}[colback=soclecol!10, colframe=soclecol, arc=4pt,
    boxrule=1pt, left=4pt, right=4pt, top=2pt, bottom=2pt, width=3.5cm]
    \centering\bfseries\color{soclecol} Socle
  \end{tcolorbox}}
\newcommand{\versionstandard}{%
  \begin{tcolorbox}[colback=stdcol!10, colframe=stdcol, arc=4pt,
    boxrule=1pt, left=4pt, right=4pt, top=2pt, bottom=2pt, width=3.5cm]
    \centering\bfseries\color{stdcol} Standard
  \end{tcolorbox}}
\newcommand{\versionappro}{%
  \begin{tcolorbox}[colback=approcol!10, colframe=approcol, arc=4pt,
    boxrule=1pt, left=4pt, right=4pt, top=2pt, bottom=2pt, width=4.5cm]
    \centering\bfseries\color{approcol} Approfondissement
  \end{tcolorbox}}

% =============================================================================
% Commandes utilitaires
% =============================================================================
\newcommand{\pts}[1]{\hfill\textit{(#1~pts)}}
\newcommand{\ptss}[1]{\hfill\textit{(#1~pt)}}
\newcommand{\reponse}{\par\noindent\dotfill\par}
\newcommand{\reponselong}{%
  \par\noindent\dotfill\par\noindent\dotfill\par}
\newcommand{\espacereponse}[1]{\vspace{#1}}

% Décimale française
\DeclareMathOperator{\e}{}

% =============================================================================
% En-tête / pied de page
% =============================================================================
\pagestyle{fancy}
\fancyhf{}
\renewcommand{\headrulewidth}{0.4pt}
\lhead{\footnotesize DS -- Ch.\,5 : Équations du 1\textsuperscript{er} degré}
\rhead{\footnotesize 2\textsuperscript{de} Bac Pro}
\cfoot{\footnotesize Page \thepage/\pageref{LastPage}}

% =============================================================================
% Titre du DS (réutilisable pour chaque version)
% =============================================================================
\newcommand{\dsheader}{%
  \begin{center}
    {\LARGE\bfseries Devoir Surveillé -- Chapitre 5}\\[6pt]
    {\large Résolution d'équations du 1\textsuperscript{er} degré
      \;\textbar\; 2\textsuperscript{de} Bac Pro}\\[10pt]
    \begin{tabular}{@{}p{5.5cm}p{5.5cm}@{}}
      \toprule
      \textbf{Durée :} 1 heure & \textbf{Calculatrice :} autorisée \\
      \textbf{Barème :} 20 points & \textbf{Documents :} non autorisés \\
      \bottomrule
    \end{tabular}\\[8pt]
    \badgeAPP{APP} \quad \badgeANA{ANA} \quad
    \badgeREA{REA} \quad \badgeVAL{VAL} \quad \badgeCOM{COM}
  \end{center}
  \vspace{4pt}
  \noindent\textbf{Nom :} \dotfill\quad
  \textbf{Prénom :} \dotfill\quad
  \textbf{Classe :} \dotfill
  \vspace{8pt}
  \hrule height 0.5pt
  \vspace{10pt}
}

% =============================================================================
\begin{document}

% ╔═══════════════════════════════════════════════════════════════════════════╗
% ║                         VERSION SOCLE                                    ║
% ╚═══════════════════════════════════════════════════════════════════════════╝
\dsheader
\begin{center}\versionsocle\end{center}
\vspace{6pt}

% ----- Exercice 1 -----
\subsection*{Exercice 1 -- Résoudre des équations (guidé) \hfill 10 points}
\textit{Détailler les étapes en suivant le modèle.}

\medskip

\textbf{1.} \badgeREA{REA} Résoudre $4x - 6 = 10$ en complétant :\pts{3}

\begin{quote}
Étape 1 : On isole $4x$ \quad $\longrightarrow$ \quad $4x = 10 + 6 = $ \dotfill

Étape 2 : On divise par 4 \quad $\longrightarrow$ \quad $x = \dfrac{\ldots}{4} = $ \dotfill
\end{quote}

\vspace{6pt}

\textbf{2.} \badgeREA{REA} Résoudre $5x + 3 = 28$ en complétant :\pts{3}

\begin{quote}
Étape 1 : $5x = 28 - $ \ldots\ $= $ \dotfill

Étape 2 : $x = \dfrac{\ldots}{\ldots} = $ \dotfill
\end{quote}

\vspace{6pt}

\textbf{3.} \badgeREA{REA} Résoudre $2x = 9$. \textit{(Le résultat n'est pas entier, donner la valeur exacte.)}\pts{2}

\reponse

\vspace{6pt}

\textbf{4.} \badgeVAL{VAL} Vérification : remplacer $x$ par ta réponse dans $4x - 6 = 10$ (question 1). Retrouves-tu bien 10 ?\pts{2}

\begin{quote}
$4 \times$ \ldots\ $- 6 = $ \ldots\ $- 6 = $ \ldots \quad $\checkmark$ ou $\times$
\end{quote}

\vspace{8pt}

% ----- Exercice 2 -----
\subsection*{Exercice 2 -- Problème guidé \hfill 10 points}

\begin{tcolorbox}[colback=blue!5, colframe=blue!40, arc=4pt, boxrule=0.6pt,
  left=6pt, right=6pt, top=4pt, bottom=4pt]
\textbf{Situation :} Un menuisier achète des panneaux de bois à \textbf{8\,\euro{} pièce}.
La \textbf{livraison} coûte \textbf{12\,\euro{}}. Il paie \textbf{60\,\euro{}} au total.
\end{tcolorbox}

\medskip

\textbf{1.} \badgeAPP{APP} Compléter l'équation qui traduit la situation :\pts{2}

\begin{quote}
\ldots\ $\times\; x \;+$ \ldots\ $=$ \ldots
\end{quote}

\vspace{4pt}

\textbf{2.} \badgeREA{REA} Résoudre cette équation.\pts{3}

\reponselong

\vspace{4pt}

\textbf{3.} \badgeCOM{COM} Écrire une phrase de réponse.\pts{2}

\reponse

\vspace{4pt}

\textbf{4.} \badgeANA{ANA} Le prix du panneau augmente à \textbf{10\,\euro{}}. Avec le même budget de 60\,\euro{}, combien de panneaux \textbf{entiers} peut-il acheter ?\pts{3}

\reponselong

% ----- Corrections Socle -----
\vspace{12pt}
\begin{correction}
\textbf{Exercice 1}

\textbf{1.} $4x = 10 + 6 = 16$ \quad puis \quad $x = \dfrac{16}{4} = 4$

\textbf{2.} $5x = 28 - 3 = 25$ \quad puis \quad $x = \dfrac{25}{5} = 5$

\textbf{3.} $x = \dfrac{9}{2} = 4{,}5$

\textbf{4.} $4 \times 4 - 6 = 16 - 6 = 10$ \quad $\checkmark$

\bigskip
\textbf{Exercice 2}

\textbf{1.} $8 \times x + 12 = 60$, soit $8x + 12 = 60$

\textbf{2.} $8x = 60 - 12 = 48$ \quad donc \quad $x = \dfrac{48}{8} = 6$

\textbf{3.} Le menuisier peut acheter \textbf{6 panneaux}.

\textbf{4.} Nouvelle équation : $10x + 12 = 60$

$10x = 48$ \quad donc \quad $x = \dfrac{48}{10} = 4{,}8$

On ne peut pas acheter $4{,}8$ panneaux : il peut acheter \textbf{4 panneaux entiers}.
\end{correction}

% ╔═══════════════════════════════════════════════════════════════════════════╗
% ║                         VERSION STANDARD                                 ║
% ╚═══════════════════════════════════════════════════════════════════════════╝
\newpage
\dsheader
\begin{center}\versionstandard\end{center}
\vspace{6pt}

% ----- Exercice 1 -----
\subsection*{Exercice 1 -- Résolution d'équations \hfill 8 points}

\textbf{1.} \badgeREA{REA} Résoudre $4x - 6 = 10$.\pts{2}

\reponselong

\vspace{4pt}

\textbf{2.} \badgeREA{REA} Résoudre $3(x + 2) = 2x + 11$.\pts{3}

\reponselong

\vspace{4pt}

\textbf{3.} \badgeREA{REA} Résoudre $\dfrac{2x + 1}{3} = 5$.\pts{3}

\reponselong

\vspace{8pt}

% ----- Exercice 2 -----
\subsection*{Exercice 2 -- Mise en équation \hfill 6 points}

\begin{tcolorbox}[colback=blue!5, colframe=blue!40, arc=4pt, boxrule=0.6pt,
  left=6pt, right=6pt, top=4pt, bottom=4pt]
\textbf{Situation :} Un artisan loue du matériel. La location coûte un \textbf{forfait de 15\,\euro{} par jour} plus \textbf{2,50\,\euro{} par outil} loué. Il dispose d'un budget de \textbf{40\,\euro{}}.
\end{tcolorbox}

\medskip

\textbf{1.} \badgeAPP{APP} Exprimer le coût $C$ en fonction du nombre d'outils $n$.\pts{2}

\reponse

\vspace{4pt}

\textbf{2.} \badgeREA{REA} Combien d'outils peut-il louer avec exactement 40\,\euro{} ? Résoudre l'équation.\pts{2}

\reponselong

\vspace{4pt}

\textbf{3.} \badgeVAL{VAL} Vérifier le résultat en remplaçant $n$ dans l'expression de $C$.\pts{2}

\reponse

\vspace{8pt}

% ----- Exercice 3 -----
\subsection*{Exercice 3 -- Problème professionnel \hfill 6 points}

\begin{tcolorbox}[colback=blue!5, colframe=blue!40, arc=4pt, boxrule=0.6pt,
  left=6pt, right=6pt, top=4pt, bottom=4pt]
\textbf{Situation :} Un menuisier découpe une planche de \textbf{2,40\,m} en \textbf{6 planches} de même longueur $x$ (en cm). Chaque coupe entraîne une \textbf{perte de 0,5\,cm} (épaisseur de la lame).
\end{tcolorbox}

\medskip

\textbf{1.} \badgeAPP{APP} Écrire l'équation traduisant la situation. \textit{(Attention : combien de coupes pour 6 planches ?)}\pts{2}

\reponse

\vspace{4pt}

\textbf{2.} \badgeREA{REA} Résoudre l'équation. Donner la longueur de chaque planche arrondie au centième.\pts{2}

\reponselong

\vspace{4pt}

\textbf{3.} \badgeANA{ANA} Si l'on arrondit la longueur à 40\,cm, est-ce réalisable ? Justifier.\pts{2}

\reponselong

% ----- Corrections Standard -----
\vspace{12pt}
\begin{correction}
\textbf{Exercice 1}

\textbf{1.} $4x - 6 = 10 \;\Longrightarrow\; 4x = 16 \;\Longrightarrow\; x = 4$

\textbf{2.} $3(x + 2) = 2x + 11 \;\Longrightarrow\; 3x + 6 = 2x + 11
  \;\Longrightarrow\; 3x - 2x = 11 - 6 \;\Longrightarrow\; x = 5$

\textbf{3.} $\dfrac{2x + 1}{3} = 5 \;\Longrightarrow\; 2x + 1 = 15
  \;\Longrightarrow\; 2x = 14 \;\Longrightarrow\; x = 7$

\bigskip
\textbf{Exercice 2}

\textbf{1.} $C = 15 + 2{,}5\,n$

\textbf{2.} $15 + 2{,}5\,n = 40 \;\Longrightarrow\; 2{,}5\,n = 25
  \;\Longrightarrow\; n = \dfrac{25}{2{,}5} = 10$

Il peut louer \textbf{10 outils}.

\textbf{3.} Vérification : $C = 15 + 2{,}5 \times 10 = 15 + 25 = 40$\,\euro{} \quad $\checkmark$

\bigskip
\textbf{Exercice 3}

\textbf{1.} 6 planches $\Rightarrow$ 5 coupes. Équation : $6x + 5 \times 0{,}5 = 240$, soit $6x + 2{,}5 = 240$.

\textbf{2.} $6x = 240 - 2{,}5 = 237{,}5 \;\Longrightarrow\; x = \dfrac{237{,}5}{6} \approx 39{,}58$\,cm.

\textbf{3.} Si $x = 40$ : $6 \times 40 + 2{,}5 = 242{,}5 > 240$. La planche est trop courte : c'est \textbf{impossible}.
\end{correction}

% ╔═══════════════════════════════════════════════════════════════════════════╗
% ║                       VERSION APPROFONDISSEMENT                          ║
% ╚═══════════════════════════════════════════════════════════════════════════╝
\newpage
\dsheader
\begin{center}\versionappro\end{center}
\vspace{6pt}

% ----- Exercice 1 -----
\subsection*{Exercice 1 -- Équations variées \hfill 8 points}

\textbf{1.} \badgeREA{REA} Résoudre $3(2x - 1) = 4(x + 3) - 5$.\pts{2}

\reponselong

\vspace{4pt}

\textbf{2.} \badgeREA{REA} Résoudre $\dfrac{5x - 2}{4} = \dfrac{x + 6}{2}$.\pts{2}

\reponselong

\vspace{4pt}

\textbf{3.} \badgeANA{ANA} Résoudre $7x - 3 = 7x + 5$. Que constate-t-on ? Conclure.\pts{2}

\reponselong

\vspace{4pt}

\textbf{4.} \badgeANA{ANA} Résoudre $2(3x + 1) = 6x + 2$. Que constate-t-on ? Conclure.\pts{2}

\reponselong

\vspace{8pt}

% ----- Exercice 2 -----
\subsection*{Exercice 2 -- Problème ouvert : comparaison de devis \hfill 12 points}

\begin{tcolorbox}[colback=blue!5, colframe=blue!40, arc=4pt, boxrule=0.6pt,
  left=6pt, right=6pt, top=4pt, bottom=4pt]
\textbf{Situation :} Un artisan menuisier compare deux fournisseurs pour l'achat de planches.

\medskip
\begin{tabular}{@{}lcc@{}}
  \toprule
  & \textbf{Fournisseur A} & \textbf{Fournisseur B} \\
  \midrule
  Prix par planche & 12\,\euro{} & 9\,\euro{} \\
  Frais de livraison & 30\,\euro{} & 60\,\euro{} \\
  \bottomrule
\end{tabular}
\end{tcolorbox}

\medskip

\textbf{1.} \badgeAPP{APP} Exprimer le coût total $C_A$ et $C_B$ en fonction du nombre $n$ de planches.\pts{2}

\reponse
\reponse

\vspace{4pt}

\textbf{2.} \badgeREA{REA} Pour quelle valeur de $n$ les deux fournisseurs proposent-ils le même prix ? À partir de combien de planches le fournisseur B devient-il moins cher ?\pts{3}

\reponselong
\reponse

\vspace{4pt}

\textbf{3.} \badgeVAL{VAL} Vérifier en calculant $C_A$ et $C_B$ pour $n = 9$ puis $n = 11$.\pts{2}

\reponselong

\vspace{4pt}

\textbf{4.} \badgeREA{REA} L'artisan dispose d'un budget de \textbf{200\,\euro{}}. Combien de planches peut-il acheter au maximum chez chaque fournisseur ?\pts{3}

\reponselong
\reponse

\vspace{4pt}

\textbf{5.} \badgeCOM{COM} Rédiger un conseil argumenté pour l'artisan.\pts{2}

\reponse
\reponselong

% ----- Corrections Approfondissement -----
\vspace{12pt}
\begin{correction}
\textbf{Exercice 1}

\textbf{1.} $3(2x - 1) = 4(x + 3) - 5$

$6x - 3 = 4x + 12 - 5 \;\Longrightarrow\; 6x - 3 = 4x + 7
  \;\Longrightarrow\; 2x = 10 \;\Longrightarrow\; x = 5$

\textbf{2.} $\dfrac{5x - 2}{4} = \dfrac{x + 6}{2}$

On multiplie par 4 : $5x - 2 = 2(x + 6) = 2x + 12$

$3x = 14 \;\Longrightarrow\; x = \dfrac{14}{3}$

\textbf{3.} $7x - 3 = 7x + 5 \;\Longrightarrow\; -3 = 5$

C'est une \textbf{contradiction} : l'équation n'a \textbf{aucune solution}.

\textbf{4.} $2(3x + 1) = 6x + 2 \;\Longrightarrow\; 6x + 2 = 6x + 2 \;\Longrightarrow\; 0 = 0$

C'est une \textbf{identité} : l'équation est vérifiée pour \textbf{tout réel $x$} (infinité de solutions).

\bigskip
\textbf{Exercice 2}

\textbf{1.} $C_A = 12n + 30$ \qquad $C_B = 9n + 60$

\textbf{2.} $12n + 30 = 9n + 60 \;\Longrightarrow\; 3n = 30 \;\Longrightarrow\; n = 10$

Les deux coûts sont égaux pour $n = 10$ planches. Le fournisseur B devient moins cher \textbf{à partir de 11 planches}.

\textbf{3.}
\begin{itemize}[nosep]
  \item $n = 9$ : $C_A = 12 \times 9 + 30 = 138$\,\euro{} \quad $C_B = 9 \times 9 + 60 = 141$\,\euro{} \quad ($A < B$)
  \item $n = 11$ : $C_A = 12 \times 11 + 30 = 162$\,\euro{} \quad $C_B = 9 \times 11 + 60 = 159$\,\euro{} \quad ($B < A$) \quad $\checkmark$
\end{itemize}

\textbf{4.} Budget de 200\,\euro{} :
\begin{itemize}[nosep]
  \item Fournisseur A : $12n + 30 \leq 200 \;\Longrightarrow\; 12n \leq 170 \;\Longrightarrow\; n \leq 14{,}1\overline{6}$ \quad $\Rightarrow$ \textbf{14 planches}
  \item Fournisseur B : $9n + 60 \leq 200 \;\Longrightarrow\; 9n \leq 140 \;\Longrightarrow\; n \leq 15{,}5\overline{5}$ \quad $\Rightarrow$ \textbf{15 planches}
\end{itemize}

\textbf{5.} Conseil : pour moins de 10 planches, choisir le fournisseur A (livraison moins chère). À partir de 10 planches, choisir le fournisseur B (prix unitaire plus avantageux). Avec un budget de 200\,\euro{}, le fournisseur B permet d'acheter \textbf{1 planche de plus} (15 contre 14).
\end{correction}

\end{document}
